% Options for packages loaded elsewhere
\PassOptionsToPackage{unicode}{hyperref}
\PassOptionsToPackage{hyphens}{url}
\documentclass[
  11pt,
  a4paper,
]{article}
\usepackage{xcolor}
\usepackage[margin=22mm]{geometry}
\usepackage{amsmath,amssymb}
\setcounter{secnumdepth}{-\maxdimen} % remove section numbering
\usepackage{iftex}
\ifPDFTeX
  \usepackage[T1]{fontenc}
  \usepackage[utf8]{inputenc}
  \usepackage{textcomp} % provide euro and other symbols
\else % if luatex or xetex
  \usepackage{unicode-math} % this also loads fontspec
  \defaultfontfeatures{Scale=MatchLowercase}
  \defaultfontfeatures[\rmfamily]{Ligatures=TeX,Scale=1}
\fi
\usepackage{lmodern}
\ifPDFTeX\else
  % xetex/luatex font selection
\fi
% Use upquote if available, for straight quotes in verbatim environments
\IfFileExists{upquote.sty}{\usepackage{upquote}}{}
\IfFileExists{microtype.sty}{% use microtype if available
  \usepackage[]{microtype}
  \UseMicrotypeSet[protrusion]{basicmath} % disable protrusion for tt fonts
}{}
\makeatletter
\@ifundefined{KOMAClassName}{% if non-KOMA class
  \IfFileExists{parskip.sty}{%
    \usepackage{parskip}
  }{% else
    \setlength{\parindent}{0pt}
    \setlength{\parskip}{6pt plus 2pt minus 1pt}}
}{% if KOMA class
  \KOMAoptions{parskip=half}}
\makeatother
\usepackage{longtable,booktabs,array}
\newcounter{none} % for unnumbered tables
\usepackage{calc} % for calculating minipage widths
% Correct order of tables after \paragraph or \subparagraph
\usepackage{etoolbox}
\makeatletter
\patchcmd\longtable{\par}{\if@noskipsec\mbox{}\fi\par}{}{}
\makeatother
% Allow footnotes in longtable head/foot
\IfFileExists{footnotehyper.sty}{\usepackage{footnotehyper}}{\usepackage{footnote}}
\makesavenoteenv{longtable}
\setlength{\emergencystretch}{3em} % prevent overfull lines
\providecommand{\tightlist}{%
  \setlength{\itemsep}{0pt}\setlength{\parskip}{0pt}}
\usepackage{bookmark}
\IfFileExists{xurl.sty}{\usepackage{xurl}}{} % add URL line breaks if available
\urlstyle{same}
\hypersetup{
  pdftitle={Effective lower bounds I: Baker\textquotesingle s theorem and its arithmetic tools},
  pdfauthor={GPT-6.1 Sol (OpenAI), Ultra},
  hidelinks,
  pdfcreator={LaTeX via pandoc}}

\title{Effective lower bounds I: Baker\textquotesingle s theorem and its
arithmetic tools}
\author{GPT-6.1 Sol (OpenAI), Ultra}
\date{October 2026}

\providecommand{\Eta}{\mathrm{H}}
\begin{document}
\maketitle

\emph{Draft. Public domain (CC0).}

An effective logarithm estimate combines a rank argument with a
comparison of arithmetic and analytic size. We begin with the rank test:
how many exact zeros would force an exponential polynomial to be zero?
We then choose polynomial bases that preserve its coefficients, prove
the denominator bounds needed to clear its derivatives, and account for
the height cost of normalizing the logarithmic form.

Theorem 4.1 states the target, including the algebraic constant and
arbitrary fixed logarithm branches. The next lesson,
\href{../TR-BAKER-05.html}{\emph{Effective lower bounds II: proof of
Baker's theorem}}, gives its induction and parameter choices. The tools
proved here include a degree-adapted polynomial family with a derivative
denominator independent of the evaluation point; this supplies an
alternative to the interval denominator in the classical construction.

We use the qualitative theorem from
\href{../TR-BAKER-03.html}{\emph{Baker's theorem II: extrapolation and
the end of the proof}}, elementary polynomial algebra and integer
factorization. The algebraic-integer criterion is proved in
\href{../../NT-ANT/algebraic-integers-and-rings-of-integers.html}{\emph{Algebraic
integers and rings of integers}}, Proposition 1.1. The freely accessible
{[}Waldschmidt 2000{]} supplies the effective logarithm estimates and
Matveev polynomial method discussed here. The logarithm methods of
Baker, Feldman and Matveev are credited in the references. All the
arithmetic arguments used below are supplied in full.

\subsection{1. The theorem, with its
quantifiers}\label{the-theorem-with-its-quantifiers}

For an algebraic number \(\gamma\), let \(H(\gamma)\) be the maximum
absolute coefficient of its primitive irreducible polynomial in
\(\mathbb Z[X]\), with positive leading coefficient. This is the
\textbf{naive height}. Thus \[
H(p/q)=\max(|p|,q)\quad(\gcd(p,q)=1,\ q>0),\qquad
H(0)=H(1)=H(i)=1.
\] It differs from the logarithmic absolute Weil height used in the
first lesson.

\textbf{Theorem 4.1 (target: Baker's effective lower bound).} Fix an
integer \(n\ge1\), an integer \(d\ge1\), nonzero algebraic numbers
\(\alpha_1,\ldots,\alpha_n\) of degree at most \(d\) and naive height at
most \(A\), where \(A\ge2\), and a determination \(\ell_j\) of each
logarithm: \[
e^{\ell_j}=\alpha_j.
\] There is an effectively computable \(C>0\), depending on \(n,d,A\)
and these determinations, with the following property. For every
\(B\ge2\) and every tuple of algebraic numbers
\(\beta_0,\ldots,\beta_n\), each of degree at most \(d\) and naive
height at most \(B\), the form \[
\Lambda=\beta_0+\beta_1\ell_1+\cdots+\beta_n\ell_n
\] satisfies \[
\Lambda=0\quad\hbox{or}\quad |\Lambda|>B^{-C}. \tag{4.1}
\] The constant is uniform in \(B\) and in all the allowed coefficients.
No rational independence assumption on the logarithms is required.

``Effective'' means that the proof gives a terminating procedure for
choosing \(C\) from the fixed algebraic data and the logarithm branches.
It does not promise that the resulting number is small. Increasing \(C\)
preserves (4.1), which is useful for absorbing fixed numerical factors.

The separation into zero and nonzero cases matters. For example,
\(2\log2-\log4=0\) on the real branches. A lower-bound theorem cannot
exclude that relation. The proof will either find a bounded integer
relation and reduce the number of logarithms, or force the proposed very
small form to be zero.

For principal logarithms, or a specified bounded family of branches, the
proof gives a power of \(\log A\): \[
C=C'(n,d)(\log A)^{\kappa(n)} \tag{4.2}
\] in the principal case, after increasing the effective constant if
necessary. The next lesson establishes this dependence. For arbitrary
branches the leading constant also depends on their branch integers.
This qualification is necessary; it is not removable by an adjustment
depending only on \(n,d\).

To see why, set both bases equal to \(1\), take \(A=B=2\), and use the
coefficients \(\sqrt2,-1\). Write \[
(1+\sqrt2)^r=p_r+q_r\sqrt2,\qquad p_r,q_r\in\mathbb Z.
\] Choose logarithms \(2\pi iq_r\) and \(2\pi ip_r\) of the two bases.
The nonzero form becomes \[
2\pi i(q_r\sqrt2-p_r)=-2\pi i(1-\sqrt2)^r,
\] whose absolute value tends to zero while the degrees and naive
heights stay fixed. Hence a common \(C\) over all branch choices is
impossible. The fixed-branch quantifier in Theorem 4.1 avoids this
problem.

\subsection{2. What the estimate gives}\label{what-the-estimate-gives}

The effective theorem is the target of the construction below; its proof
is completed in the next lesson. The resulting transcendence measures
for logarithms and the fixed-base multiplicative bound are proved after
that argument, in
\href{../TR-BAKER-05.html\#10-consequences-and-later-refinements}{\emph{Effective
lower bounds II: proof of Baker's theorem}}, Section 10.

Later applications need sharper height dependence than a single power of
\(\log A\). The individual-height product \(\prod_j\log A_j\), and the
different coefficient and branch hypotheses of the modern estimates, are
treated in {[}Waldschmidt 2000{]} and {[}Waldschmidt 2003{]}. The later
lesson \href{../TR-BAKER-08.html}{\emph{Linear forms in many logarithms:
the modern estimates and how to use them}} gives the precise modern
estimates and their proofs. The present theorem first establishes the
decisive polynomial dependence on the coefficient height.

\subsection{3. Plan the final rank test}\label{plan-the-final-rank-test}

Before estimating an auxiliary function, decide what enough zeros would
imply. Suppose its values have the form \[
F(i)=\sum_{s=0}^{l-1}P_s(i)w_s^i,\qquad \deg P_s<k.
\] There are \(kl\) polynomial coefficients. Evaluation at
\(i=0,\ldots,kl-1\) is a square linear map from those coefficients to
the values of \(F\). If its matrix is invertible, vanishing of all those
values forces every coefficient to vanish. In a construction with a
nonzero coefficient vector, that is the contradiction we want.

The two possible failures have different meanings. Coincident nodes
combine frequency blocks, so the values can no longer distinguish their
coefficients. When \(k\ge2\), a zero node destroys the polynomially
weighted columns. The next formula identifies both failures exactly. It
also gives the complete rank test when each frequency carries several
polynomial coefficients.

\textbf{Lemma 4.6.} Let \(k,l\ge1\), and let \(w_0,\ldots,w_{l-1}\) be
distinct nonzero complex numbers. Order the columns by pairs \((s,r)\),
first \(s=0,\ldots,l-1\), then \(r=0,\ldots,k-1\). The matrix with rows
\(i=0,\ldots,kl-1\) and entries \[
A_{i,(s,r)}=i^r w_s^i,\qquad 0^0=1,
\] has determinant \[
\det A=
\left(\prod_{s=0}^{l-1}w_s^{k(k-1)/2}\right)
\left(\prod_{r=0}^{k-1}r!\right)^l
\prod_{0\le s<t<l}(w_t-w_s)^{k^2}. \tag{4.15}
\] In particular, it is nonzero.

\textbf{Proof.} Put \(N=kl\) and \(K_0=k(k-1)/2\). The ordinary
Vandermonde determinant at nodes \(z_0,\ldots,z_{N-1}\) is
\(\prod_{a<b}(z_b-z_a)\). Indeed, its polynomial determinant changes
sign under exchanging nodes, is divisible by each difference, and has
the same total degree as that product; the coefficient obtained by
choosing the successive highest powers is \(1\).

Apply this identity to the nodes \[
z_{s,r}=w_s+r\varepsilon\quad(0\le s<l,\ 0\le r<k)
\] in our column order, and compare the coefficient of the lowest
nonzero power \(\varepsilon^{lK_0}\). Taylor expansion of the columns
shows that in each block the derivative orders must be distinct, since
equal orders give repeated columns. The smallest total order is \(K_0\),
attained exactly by \(0,\ldots,k-1\). Its coefficient within that block
is \[
\det\left(\frac{r^t}{t!}\right)_{0\le t,r<k}
=\frac{\prod_{a<b}(b-a)}{\prod_{t=0}^{k-1}t!}=1.
\] Thus the coefficient on the determinant side is the determinant of
the confluent matrix \[
F_{i,(s,r)}=(i)_r w_s^{i-r},\qquad
(i)_r=i(i-1)\cdots(i-r+1),
\] with \((i)_0=1\). On the product side, the within-block differences
contribute \(\left(\prod_{r=0}^{k-1}r!\right)^l\varepsilon^{lK_0}\);
each pair of different blocks contributes \((w_t-w_s)^{k^2}\).
Consequently \[
\det F=\left(\prod_{r=0}^{k-1}r!\right)^l
\prod_{s<t}(w_t-w_s)^{k^2}. \tag{4.16}
\]

The falling factorials \((i)_0,\ldots,(i)_r\) form a triangular monic
basis for polynomials in \(i\) of degree at most \(r\). Write \[
i^r=\sum_{t=0}^r S(r,t)(i)_t,\qquad S(r,r)=1.
\] Then the \(r\)-th column of the \(s\)-block of \(A\) is
\(\sum_{t\le r}S(r,t)w_s^t F_{(s,t)}\). This triangular column change
has determinant \(\prod_{r=0}^{k-1}w_s^r=w_s^{K_0}\) for each block.
Multiplying (4.16) by these factors proves (4.15). \(\square\)

For \(k=l=2\), the exact formula is \[
\det\begin{pmatrix}
1&0&1&0\\
w_0&w_0&w_1&w_1\\
w_0^2&2w_0^2&w_1^2&2w_1^2\\
w_0^3&3w_0^3&w_1^3&3w_1^3
\end{pmatrix}
=w_0w_1(w_1-w_0)^4. \tag{4.17}
\] At \(w_0=2,w_1=3\), the determinant is \(6\). Nonzero nodes are
needed: when \(k\ge2\), a node \(w_s=0\) makes the column with \(r=1\)
identically zero.

\subsubsection{Recover four coefficients from four
values}\label{recover-four-coefficients-from-four-values}

Take \[
F(i)=(a+bi)2^i+(c+di)3^i,\qquad y_j=F(j).
\] The determinant \(6\) calculated above proves uniqueness. Solving the
same system gives the more informative identities \[
\begin{aligned}
a&=-27y_0+36y_1-15y_2+2y_3,\\
b&=-9y_0+\tfrac{21}{2}y_1-4y_2+\tfrac12y_3,\\
c&=28y_0-36y_1+15y_2-2y_3,\\
d&=-4y_0+\tfrac{16}{3}y_1-\tfrac73y_2+\tfrac13y_3.
\end{aligned}
\] For example, \((y_0,y_1,y_2,y_3)=(3,9,32,119)\) recovers
\((a,b,c,d)=(1,-1,2,1)\). These identities can be verified by inserting
the four rows in (4.17); they give an explicit inverse, rather than a
numerical rank calculation.

The count of values is essential. With only \(y_0=y_1=y_2=0\), the
nonzero choice \((a,b,c,d)=(6,\tfrac32,-6,1)\) is possible and gives
\(y_3=3\). Thus three zeros do not yet force the four coefficients to
vanish. In the effective proof, extrapolation must reach the number of
rows required by the rank test. A small analytic estimate before that
point would not finish the argument.

The polynomial coefficient blocks need not initially be written in
monomials. We next construct a basis that retains the same coefficient
count while accommodating the integer-valued polynomials used for
denominator control.

\subsection{4. Preserve the polynomial
coefficients}\label{preserve-the-polynomial-coefficients}

\textbf{Lemma 4.5.} Let \(K\) be a field of characteristic zero,
\(P\in K[X]\) a polynomial of degree \(n>0\), and \(0\le m\le n\). The
\(n+1\) polynomials \[
P(X),P(X+1),\ldots,P(X+m),\quad
1,X,\ldots,X^{n-m-1} \tag{4.14}
\] are linearly independent over \(K\). If \(m=n\), the second list is
empty. They consequently form a basis of \(K[X]_{\le n}\).

\textbf{Proof.} Suppose \(\sum_{j=0}^m A_jP(X+j)\) has degree at most
\(n-m-1\). For \(0\le t\le m\), the coefficient of \(X^{n-t}\) in
\(P(X+j)\), as a polynomial \(q_t(j)\), has degree \(t\) and leading
coefficient \(a_n\binom nt\ne0\). Therefore \(q_0,\ldots,q_m\) form a
basis of the polynomials of degree at most \(m\) in \(j\).

Vanishing of those highest \(m+1\) coefficients gives
\(\sum_j A_jq_t(j)=0\) for all \(t\). By the basis assertion it gives
\(\sum_j A_jj^t=0\) for \(t=0,\ldots,m\). The Vandermonde matrix at the
distinct elements \(0,\ldots,m\) is invertible, so every \(A_j=0\). The
remaining monomials are independent. \(\square\)

For \(P(X)=X^2,\ m=1\), the basis is \(X^2,(X+1)^2,1\). Subtracting the
first and third elements from the second produces \(2X\); characteristic
zero permits division by \(2\).

The characteristic assumption is essential. In characteristic \(p\),
\(P(X)=X^p-X\) satisfies \(P(X+1)=P(X)\), so even two of the proposed
translates coincide. All coefficient fields in the effective proof are
subfields of \(\mathbb C\), where Lemma 4.5 applies. Exercise 3 gives
the complete induction proof by finite differences.

\subsubsection{Read the basis as a finite-difference
calculation}\label{read-the-basis-as-a-finite-difference-calculation}

For \(P(X)=X^3\) and \(m=2\), the proposed basis is \[
X^3,\quad (X+1)^3,\quad (X+2)^3,\quad 1.
\] Its first difference is \((X+1)^3-X^3=3X^2+3X+1\). Its second
difference is \[
(X+2)^3-2(X+1)^3+X^3=6X+6.
\] Together with \(1\), these recover \(X\), then \(X^2\), and finally
\(X^3\). Each difference removes one highest degree. Lemma 4.5
formalizes that mechanism for any polynomial with nonzero leading
coefficient; Solution 3 gives the induction in full.

This explains a practical division of labour. The translate basis
preserves all polynomial coefficients needed by the rank test. The
choice of its underlying polynomial controls the arithmetic size of
those coefficients and their derivatives. We now choose that polynomial.

\subsection{5. Form normalized
derivatives}\label{form-normalized-derivatives}

For the auxiliary construction, derivatives must satisfy two
requirements at once: their values must be small enough analytically,
and their rational denominators must be cheap enough to clear. The
normalization by \(m!\) makes the first requirement a coefficient
estimate. A product of consecutive factors makes the second requirement
accessible to divisibility arguments.

After Lemma 4.3, Section 6 will give a stronger common denominator
independent of the integer evaluation point. We first prove the direct
interval bound because its positive coefficient expansion shows exactly
how the normalized derivative is formed, and because the effective
argument can use it without changing its original parameter choices.

For the classical construction, use \[
\Delta(x;k)=\frac{(x+1)\cdots(x+k)}{k!},\qquad
\Delta(x;k,l,m)=\frac1{m!}\frac{d^m}{dx^m}\Delta(x;k)^l, \tag{4.7}
\] where \(k\ge1\) and \(l,m\ge0\) are integers. For a nonnegative
integer \(x\), set \[
\nu(x;k)=\operatorname{lcm}(x+1,\ldots,x+k).
\] The normalized derivative is exactly the coefficient of \(t^m\) in
\(\Delta(x+t;k)^l\).

\textbf{Lemma 4.3.} For every integer \(x\ge0\), \[
\nu(x;k)^m\Delta(x;k,l,m)\in\mathbb Z_{\ge0},\qquad
0\le\Delta(x;k,l,m)\le4^{l(x+k)}. \tag{4.8}
\] The integer is positive when \(0\le m\le kl\), and is zero when
\(m>kl\).

\textbf{Proof.} First \(\Delta(x;k)=\binom{x+k}{k}\) is a positive
integer. Regard the list \(x+1,\ldots,x+k\) as repeated \(l\) times.
Expanding \[
\Delta(x+t;k)^l
=\Delta(x;k)^l\prod_{j=1}^k\left(1+\frac{t}{x+j}\right)^l
\] gives \[
\Delta(x;k,l,m)=\Delta(x;k)^l
\sum_{\substack{\text{choices of }m\\\text{positions in the repeated list}}}
\frac1{(x+j_1)\cdots(x+j_m)}. \tag{4.9}
\] There are \(\binom{kl}{m}\) choices, with the sum read as zero when
\(m>kl\). Multiplication by \(\nu(x;k)^m\) clears every denominator,
proving integrality and the positivity assertion. For \(l=0\), the only
nonzero derivative is \(m=0\), with value \(1\).

Since \(x+j\ge1\), each reciprocal product is at most \(1\). The
inequalities \(\binom{x+k}{k}\le2^{x+k}\) and \(\binom{kl}{m}\le2^{kl}\)
therefore give \[
\Delta(x;k,l,m)\le2^{l(x+k)+kl}\le4^{l(x+k)}.
\] \(\square\)

For example, at \(k=l=2\), \[
\Delta(x;2)^2=\frac{x^4+6x^3+13x^2+12x+4}{4}.
\] At \(x=1\), its normalized derivatives for \(m=0,\ldots,5\) are \[
9,\quad15,\quad\frac{37}{4},\quad\frac52,\quad\frac14,\quad0.
\] Here \(\nu(1;2)=6\), so multiplying the \(m\)-th entry by \(6^m\)
gives an integer. The final zero illustrates why an assertion of
positivity for \emph{all} orders would be false.

\subsection{6. Compare denominator
costs}\label{compare-denominator-costs}

\textbf{Lemma 4.4.} For all integers \(x\ge0,\ k\ge1\), \[
\nu(x;k)\le\left(\frac{16(x+k)}k\right)^{2k}. \tag{4.10}
\]

\textbf{Proof.} Write \(V(N)=\operatorname{lcm}(1,\ldots,N)\). We first
prove \[
V(N)\le8^N\quad(N\ge1). \tag{4.11}
\] For a prime power \(p^a\), the factorial valuation formula is \[
v_p(N!)=\sum_{j\ge1}\left\lfloor\frac{N}{p^j}\right\rfloor,
\] because multiples of \(p^j\) contribute one extra factor of \(p\).
Put \(u=\lfloor N/2\rfloor,\ v=\lceil N/2\rceil\). If the highest power
\(p^a\le N\) is at most \(v\), its full contribution to \(V(N)\) already
divides \(V(v)\). Otherwise \(p^a>v\), but \(p^{a-1}\le N/2\le v\). The
term \(j=a\) in \(v_p\binom Nu\) is \(1-0-0=1\), and every other floor
difference is nonnegative. Hence \[
V(N)\mid \binom Nu V(v).
\] Using \(\binom Nu\le2^N\), induction gives \[
V(N)\le2^N8^{\lceil N/2\rceil}\le8^N;
\] the exponent inequality \(N+3\lceil N/2\rceil\le3N\) holds for
\(N\ge2\), including \(N=3\). The base \(N=1\) is immediate.

Now set \(N=x+k\) and split \(\nu(x;k)=\nu'\nu''\), according as the
prime divisor is at most \(k\) or greater than \(k\). For \(p\le k\),
let \(a_0=\lfloor\log_p k\rfloor\). Every interval of \(k\) consecutive
integers contains a multiple of \(p^{a_0}\), so \(V(k)\mid\nu'\). If
\(p^a\) is the highest power of \(p\) dividing an integer in our
interval, then \(p^a\le N\) and \(p^{a_0}>k/p\). Thus \[
p^{a-a_0}<\frac{pN}{k}
\] whenever an extra power occurs; the weak inequality with this upper
bound also covers \(a=a_0\). Therefore \[
\nu'\le V(k)\prod_{p\le k}p\,(N/k)^{\pi(k)}
\le64^k(N/k)^k, \tag{4.12}
\] where \(\prod_{p\le k}p\mid V(k)\), \(\pi(k)\le k\), \(N/k\ge1\), and
(4.11) were used.

A prime \(p>k\) divides at most one of the consecutive integers
\(x+1,\ldots,x+k\). Its valuation in their product is consequently its
valuation in \(\nu''\), while \(v_p(k!)=0\). Hence
\(\nu''\mid\binom Nk\). Finally \[
\binom Nk\le\frac{N^k}{k!}\le(eN/k)^k. \tag{4.13}
\] For the last inequality,
\(\sum_{j=1}^k\log j\ge\int_1^k\log t\,dt=k\log k-k+1\) suffices, also
at \(k=1\). Combining (4.12)--(4.13) gives \[
\nu(x;k)\le(64e)^k(N/k)^{2k}
\le(16N/k)^{2k},
\] since \(e<4\). \(\square\)

The point of (4.10) is its dependence on the ratio \(x/k\). If \(x\) is
comparable with \(k\), then \(\log\nu(x;k)=O(k)\). A crude factorial
denominator would have logarithm \(O(k\log k)\); losing that extra
factor would spoil the effective auxiliary-function count.

Here are three different parameter regimes in the original bound,
keeping their dependence visible:

{\def\LTcaptype{none} % do not increment counter
\begin{longtable}[]{@{}
  >{\raggedright\arraybackslash}p{(\linewidth - 4\tabcolsep) * \real{0.3333}}
  >{\raggedright\arraybackslash}p{(\linewidth - 4\tabcolsep) * \real{0.3333}}
  >{\raggedright\arraybackslash}p{(\linewidth - 4\tabcolsep) * \real{0.3333}}@{}}
\toprule\noalign{}
\begin{minipage}[b]{\linewidth}\raggedright
Evaluation point
\end{minipage} & \begin{minipage}[b]{\linewidth}\raggedright
Consequence of (4.10)
\end{minipage} & \begin{minipage}[b]{\linewidth}\raggedright
Dependence on \(k\)
\end{minipage} \\
\midrule\noalign{}
\endhead
\bottomrule\noalign{}
\endlastfoot
\(x=0\) & \(\log\nu\le2k\log16\), sharpened by (4.11) to \(k\log8\) &
linear \\
\(x=k\) & \(\log\nu\le2k\log32\) & linear \\
\(x=k^2\) & \(\log\nu\le2k\log(16(k+1))\) & \(O(k\log k)\) \\
\end{longtable}
}

The last row prevents a mistaken uniform \(O(k)\) assertion when the
evaluation interval moves far from zero. Lemma 4.7 below supplies a
uniform denominator for normalized derivatives even in that regime; it
does not assert that the interval LCM itself becomes small.

\subsubsection{A denominator independent of the evaluation
point}\label{a-denominator-independent-of-the-evaluation-point}

The preceding interval estimate grows with \(x/k\). A different
divisibility argument removes this dependence from the derivative
denominator. It also accommodates arbitrary polynomial degrees. This is
the polynomial family introduced by Matveev and treated in
{[}Waldschmidt 2000{]}.

For integers \(a\ge0\) and \(b\ge1\), write \(a=bq+r\), with
\(0\le r<b\), and define \[
D_b(Z;a)=
\left(\frac{Z(Z+1)\cdots(Z+b-1)}{b!}\right)^q
\frac{Z(Z+1)\cdots(Z+r-1)}{r!}.
\] An empty product is \(1\). This polynomial has degree exactly \(a\).
In particular, for fixed \(b\), the polynomials
\(D_b(Z;0),\ldots,D_b(Z;A)\) form a basis of \(\mathbb Q[Z]_{\le A}\):
their distinct degrees and nonzero leading coefficients make the
coefficient matrix triangular and invertible.

\textbf{Lemma 4.7 (uniform derivative denominators and analytic size).}
Let \(a,C\ge0\), \(b\ge1\) be integers, and put
\(V(b)=\operatorname{lcm}(1,\ldots,b)\). For every integer \(u\),
including negative integers, and every \(0\le c\le C\), \[
\frac{V(b)^C}{c!}D_b^{(c)}(u;a)\in\mathbb Z.
\] For every \(z\in\mathbb C\), \[
\sum_{c=0}^{C}\binom Cc\left|D_b^{(c)}(z;a)\right|
\le C!e^{a+b}\left(1+\frac{|z|}{b}\right)^a.
\] Derivatives of order greater than \(a\) are zero.

\textbf{Proof.} Each consecutive-factor quotient is integer-valued on
\(\mathbb Z\). For \(u\ge1\) it is \(\binom{u+t-1}{t}\), and for
\(u\le0\) it is \((-1)^t\binom{-u}{t}\), for a block of length \(t\).
These formulas include its zeros.

To prove the sharper derivative assertion, expand the coefficient of
\(T^c\) in \(D_b(u+T;a)\). Every term is obtained by deleting \(c\)
factors from the numerator's \(a\) linear factors, then dividing the
surviving product by \((b!)^qr!\). No extra combinatorial denominator
appears, since that coefficient is \(D_b^{(c)}(u;a)/c!\).

Fix a prime \(p\), and let \(L=\lfloor\log_p b\rfloor\). For
\(1\le j\le L\), each block of \(b\) consecutive integers contains at
least \(\lfloor b/p^j\rfloor\) multiples of \(p^j\), and the shorter
block contains at least \(\lfloor r/p^j\rfloor\). Set \[
t_j=q\lfloor b/p^j\rfloor+\lfloor r/p^j\rfloor.
\] Deleting \(c\) factors leaves at least \(\max(t_j-c,0)\) such
multiples. A surviving product that is zero contributes the integer
zero. Every nonzero surviving product has \(p\)-valuation at least
\(\sum_{j=1}^{L}\max(t_j-c,0)\). The denominator has valuation
\(\sum_{j=1}^{L}t_j\), while \(V(b)^C\) has valuation \(CL\). Therefore
its cleared term has valuation at least \[
\sum_{j=1}^{L}\bigl(C+\max(t_j-c,0)-t_j\bigr)\ge0,
\] because \(c\le C\). For \(p>b\), the denominator has valuation zero.
Every term is consequently an integer after multiplication by
\(V(b)^C\), and so is their sum. This proof uses intervals of integers
and remains valid when \(u\) is negative.

For the analytic bound, set \(W=|z|+b-1\). Expanding the differentiated
product gives \[
|D_b^{(c)}(z;a)|\le
\frac{c!\binom ac W^{a-c}}{(b!)^qr!}\quad(0\le c\le a).
\] If \(W=0\), the expression is interpreted by the underlying
polynomial: a term with exponent zero has value \(1\), including at
zero. Since \(\binom Cc c!=C!/(C-c)!\le C!\), summing and then
completing the binomial sum gives \[
\sum_{c=0}^{C}\binom Cc|D_b^{(c)}(z;a)|
\le \frac{C!(W+1)^a}{(b!)^qr!}.
\] The elementary factorial inequality \(t!\ge(t/e)^t\), proved in
(4.13), yields \[
\frac1{(b!)^qr!}\le
\frac{e^a}{b^a}\left(\frac br\right)^r
\le \frac{e^{a+b}}{b^a}.
\] At \(r=0\) the middle factor is read as \(1\). Otherwise
\(r\log(b/r)\le b/e<b\), by maximizing \(t\log(1/t)\) on \(0<t\le1\).
Substitution of \(W+1=|z|+b\) proves the stated bound. The degree gives
the final zero assertion. \(\square\)

For the original polynomial, \(D_k(x+1;kl)=\Delta(x;k)^l\). Taking
\(C=m\) gives the useful specialization \[
V(k)^m\Delta(x;k,l,m)\in\mathbb Z\qquad(x\in\mathbb Z).
\] At \(k=l=2,x=1\), multiplication of the five nonzero normalized
derivatives by \(2^m\) gives \[
9,\quad30,\quad37,\quad20,\quad4.
\] The interval denominator in Lemma 4.3 used \(6^m\) at this point. At
\(x=-1\), the normalized derivatives instead are \[
0,\quad0,\quad\tfrac14,\quad\tfrac12,\quad\tfrac14;
\] the same \(2^m\) still clears them. Positivity from Lemma 4.3 is
confined to its nonnegative evaluation points; the integer-valued
assertion in Lemma 4.7 has no such restriction.

Together with (4.11), this gives a uniform logarithmic denominator cost
at most \(km\log8\). The general family lets us keep the polynomial
degree \(a\) and the block size \(b\) as separate parameters. A large
degree need not force a factorial-size derivative denominator. The
analytic bound records the corresponding size cost explicitly; it is
this pair of bounds that makes the alternative basis useful.

\subsection{7. Budget the coefficient
heights}\label{budget-the-coefficient-heights}

The rank and denominator arguments leave one bookkeeping question: how
much does division by an algebraic coefficient change the degree and
height? We first prove the elementary sum and product estimate, then
compute its effect on a lower-bound exponent.

\textbf{Lemma 4.2.} If \(\alpha,\beta\) have degree at most \(d\) and
naive height at most \(H\), where \(H\ge2\), then each of
\(\alpha+\beta,\alpha\beta\) has degree at most \(d^2\), and \[
H(\alpha+\beta),\,H(\alpha\beta)
\le (1+d^2)^{d^2}H^{4d^2}. \tag{4.6}
\] In particular, \(\log H(\alpha+\beta)\) and \(\log H(\alpha\beta)\)
are at most \(c(d)\log H\), with \[
c(d)=d^2\left(4+\frac{\log(1+d^2)}{\log2}\right).
\] For \(H\ge1\), the same logarithmic conclusion holds with
\(\log(2H)\).

\textbf{Proof.} Suppose \(aX^s+\cdots+a_s\) is the primitive minimal
polynomial of \(\alpha\). Here \(1\le a\le H\) and \(s\le d\). If
\(|\alpha|\ge1\), its defining equation gives \[
a|\alpha|^s\le H\sum_{j=0}^{s-1}|\alpha|^j\le sH|\alpha|^{s-1}.
\] Thus every conjugate of \(\alpha\) has size at most \(dH\); the same
holds for \(\beta\). Also \(a\alpha\) is an algebraic integer:
substituting \(X/a\) into the defining polynomial and multiplying by
\(a^{s-1}\) gives a monic integral polynomial in \(X\).

Let \(b\le H\) be the leading coefficient for \(\beta\), and set
\(q=ab\). Both \[
q\alpha\beta=(a\alpha)(b\beta),\qquad
q(\alpha+\beta)=b(a\alpha)+a(b\beta)
\] are algebraic integers. The compositum \(\mathbb Q(\alpha,\beta)\) is
spanned by the products of powers \(1,\alpha,\ldots,\alpha^{s-1}\) and
\(1,\beta,\ldots,\beta^{t-1}\); consequently either new number
\(\gamma\) has degree \(D\le st\le d^2\).

The monic minimal polynomial of \(q\gamma\) has integral coefficients.
Evaluating it at \(qX\) gives an integral polynomial with root
\(\gamma\) and leading coefficient \(q^D\). By Gauss's lemma, the
positive leading coefficient \(c\) of the primitive minimal polynomial
of \(\gamma\) divides \(q^D\). Hence \(c\le H^{2D}\).

Every conjugate of \(\gamma=\alpha\beta\) has size at most \(d^2H^2\). A
conjugate of \(\alpha+\beta\) has size at most \(2dH\le d^2H^2\), since
\(H\ge2\). Expanding \[
c\prod_{j=1}^D(X-\gamma^{(j)})
\] shows that each coefficient is at most \[
c\prod_j(1+|\gamma^{(j)}|)
\le H^{2D}(1+d^2H^2)^D
\le (1+d^2)^{d^2}H^{4d^2}.
\] This proves (4.6). Taking logarithms and using \(\log H\ge\log2\)
gives the displayed \(c(d)\). For \(1\le H<2\), apply the result with
\(2H\). \(\square\)

The restriction or padding at height one is real: \(H(1)=1\), but
\(H(1+1)=2\). A statement \(\log H'\le c(d)\log H\) without padding
cannot hold at \(H=1\). Inversion is particularly simple: reversing the
minimal polynomial gives \(H(\alpha^{-1})=H(\alpha)\) for
\(\alpha\ne0\), with unchanged degree. Combining this fact with Lemma
4.2 handles every finite sum, product and quotient appearing in the
later induction.

\subsubsection{Account for the constant when dividing by a
coefficient}\label{account-for-the-constant-when-dividing-by-a-coefficient}

The effective proof normalizes a nonzero logarithm coefficient to
\(-1\). A degree or height estimate alone does not finish that step:
after obtaining a lower bound for the normalized form, we must multiply
back by the coefficient and absorb every fixed factor into an exponent
of \(B\).

\textbf{Lemma 4.8 (normalization budget).} Suppose
\(\beta_0,\ldots,\beta_n\) have degree at most \(d\), naive height at
most \(B\ge2\), and \(\beta_n\ne0\). Put \[
r=4d^2,\qquad K_d=(1+d^2)^{d^2},\qquad
\gamma_j=-\beta_j/\beta_n\quad(0\le j<n).
\] Each \(\gamma_j\) has degree at most \(d^2\) and height at most
\(K_dB^r\). For fixed logarithms define \[
\Lambda_* =\gamma_0+\sum_{j=1}^{n-1}\gamma_j\ell_j-\ell_n
=-\Lambda/\beta_n.
\] If \(|\Lambda_*|>(K_dB^r)^{-C_*}\), with \(C_*>0\), then \[
|\Lambda|>B^{-C},\qquad
C=1+rC_*+\frac{\log d+C_*\log K_d}{\log2}.
\]

\textbf{Proof.} Inversion preserves naive height and degree, and
changing sign preserves them too. Apply Lemma 4.2 to \(-\beta_j\) and
\(\beta_n^{-1}\). This proves the asserted degree and height bounds,
including zero coefficients.

The root bound proved in Lemma 4.2, now applied to \(\beta_n^{-1}\),
gives \(|\beta_n^{-1}|\le dB\), hence \(|\beta_n|\ge(dB)^{-1}\).
Consequently \[
|\Lambda|=|\beta_n||\Lambda_*|
>d^{-1}K_d^{-C_*}B^{-(1+rC_*)}.
\] For a fixed \(M=dK_d^{C_*}\ge1\) and \(B\ge2\), we have
\(M\le B^{\log M/\log2}\). Absorbing precisely this factor gives the
displayed exponent. \(\square\)

If every logarithm coefficient is zero, the form is simply \(\beta_0\).
It is either zero, or the same inverse root bound gives \[
|\beta_0|\ge(dB)^{-1}\ge B^{-1-\log d/\log2}.
\] Increasing the exponent slightly makes the final inequality strict,
as in Theorem 4.1. This also covers the constant-only case without
dividing by a nonexistent logarithm coefficient.

The normalization budget keeps the coefficient-height variable separate
from the fixed bases and branches. The theorem's exponent may
deteriorate when we normalize or eliminate a coefficient, but it remains
independent of \(B\). In the next lesson the analytic gain must exceed
all the denominator and height costs proved here.

\subsection{8. Exercises}\label{exercises}

\begin{enumerate}
\def\labelenumi{\arabic{enumi}.}
\item
  \textbf{Easy.} Prove directly that \(\Delta(x;k)=\binom{x+k}{k}\) is
  an integer for \(x\ge0\). Explain the cases \(x=0\) and \(k=1\).
\item
  \textbf{Medium.} Prove (4.10) with an explicit absolute constant,
  without the prime number theorem. State where you use that the
  interval contains exactly \(k\) consecutive integers.
\item
  \textbf{Medium.} Prove Lemma 4.5 by induction on \(\deg P\), using
  \(Q(X)=P(X)-P(X+1)\). Treat \(m=0\) and \(m=n\) explicitly.
\item
  \textbf{Hard.} Compute (4.17) by column operations. Explain separately
  the factors \(w_0w_1\) and \((w_1-w_0)^4\), and verify the numerical
  case \(w_0=2,w_1=3\).
\item
  \textbf{Medium.} At \(k=l=2\), compare the interval and uniform
  denominators for the normalized derivatives at \(x=1\). Compute the
  derivatives at \(x=-1\), and explain why the uniform integer-valued
  assertion survives although positivity no longer holds at every order.
\item
  \textbf{Hard.} Suppose all coefficients in Lemma 4.8 are rational of
  naive height at most \(B\). Prove that each normalized coefficient has
  height at most \(B^2\), and show that a bound
  \(|\Lambda_*|>(B^2)^{-C_*}\) gives \(|\Lambda|>B^{-(1+2C_*)}\).
  Compare this exponent with the general bound at \(d=1\).
\end{enumerate}

\subsection{9. Solutions}\label{solutions}

\textbf{Solution 1.} The product in (4.7) is \((x+k)!/x!\); dividing by
\(k!\) gives the binomial coefficient counting the \(k\)-element subsets
of a set of size \(x+k\). Hence it is integral. At \(x=0\) it equals
\(1\), and at \(k=1\) it equals \(x+1\). This also proves its
positivity.

\textbf{Solution 2.} The factorial valuation argument in Section 6
proves \(V(N)\mid\binom N{\lfloor N/2\rfloor}V(\lceil N/2\rceil)\),
hence \(V(N)\le8^N\). Split the interval LCM into primes \(p\le k\) and
\(p>k\). An interval of length \(k\) contains a multiple of every prime
power at most \(k\), so \(V(k)\) divides its small-prime part. Any extra
power is at most \(pN/k\), where \(N=x+k\); thus that part is at most
\(64^k(N/k)^k\). A prime \(p>k\) has at most one multiple in the
interval, so the large-prime part divides the product of the interval
divided by \(k!\), namely \(\binom Nk\le(eN/k)^k\). Multiplication gives
\((64e)^k(N/k)^{2k}\le(16N/k)^{2k}\). Both uses of the interval length
are essential. This proves the assigned bound with \(c=16\), including
\(k=1\).

\textbf{Solution 3.} For \(m=0\), \(P\) and the lower-degree monomials
are independent by their leading coefficients. For degree \(1\) and
\(m=1\), the difference \(P(X)-P(X+1)\) is a nonzero constant, so the
two translates are independent.

Assume the result for degree \(n-1\), and take \(1\le m\le n\). Suppose
\[
R(X)=\sum_{j=0}^m A_jP(X+j)
\] has degree at most \(n-m-1\), where a bound of \(-1\) means the zero
polynomial. Its leading coefficient gives \(\sum_jA_j=0\). Put
\(C_j=A_0+\cdots+A_j\). Telescoping gives the exact identity \[
R(X)=\left(\sum_{j=0}^mA_j\right)P(X+m)
+\sum_{j=0}^{m-1}C_j\bigl(P(X+j)-P(X+j+1)\bigr).
\] Since \(Q(X)=P(X)-P(X+1)\) has degree \(n-1\) in characteristic zero,
the induction hypothesis with parameter \(m-1\) says that its \(m\)
translates, together with the monomials up to degree \(n-m-1\), are
independent. Therefore every \(C_j=0\). Successive differences give
\(A_0=\cdots=A_{m-1}=0\), and the vanishing total gives \(A_m=0\). When
\(m=n\), the induction uses no lower-degree monomials and the same
argument works.

\textbf{Solution 4.} Write \(a=w_0,\ b=w_1\), and factor \(a\) from
column 2 and \(b\) from column 4. The remaining matrix has columns \[
C(a),\ C'(a),\ C(b),\ C'(b),\qquad
C(z)=(1,z,z^2,z^3)^{\mathsf T}.
\] Set \(\delta=b-a\), and use \[
C(b)=C(a)+\delta C'(a)+\frac{\delta^2}{2}C''(a)
+\frac{\delta^3}{6}C'''(a),
\] \[
C'(b)=C'(a)+\delta C''(a)+\frac{\delta^2}{2}C'''(a).
\] Subtract \(C(a)+\delta C'(a)\) from the third column and \(C'(a)\)
from the fourth. In the last two columns, the coefficient determinant
relative to \(C''(a),C'''(a)\) is \[
\det\begin{pmatrix}\delta^2/2&\delta\\
\delta^3/6&\delta^2/2\end{pmatrix}
=\delta^4/12.
\] Meanwhile \(\det[C(a),C'(a),C''(a),C'''(a)]=1\cdot1\cdot2\cdot6=12\),
since its matrix is triangular. The remaining determinant is
\(\delta^4\), and restoring the factors gives \(ab(b-a)^4\). For
\(a=2,b=3\), the original rows are \[
(1,0,1,0),\ (2,2,3,3),\ (4,8,9,18),\ (8,24,27,81),
\] and the formula gives \(6\).

\textbf{Solution 5.} At \(x=1\) the normalized derivatives are
\(9,15,37/4,5/2,1/4\), followed by zeros. The interval LCM is \(6\); its
cleared values are integers by Lemma 4.3. The uniform LCM is \(V(2)=2\),
and the cleared values are \(9,30,37,20,4\). Thus the same derivatives
admit the smaller denominator \(2^m\). Write \(x=-1+t\). Then \[
\Delta(-1+t;2)^2=\tfrac14t^2(1+t)^2
=\tfrac14t^2+\tfrac12t^3+\tfrac14t^4.
\] Its normalized derivatives at \(-1\) are the coefficients
\(0,0,1/4,1/2,1/4\), so multiplication by \(2^m\) gives \(0,0,1,4,4\).
The first two entries vanish even though their orders are at most the
polynomial degree. Lemma 4.3 did not assert positivity at negative
points; Lemma 4.7 proves integer-valued cleared derivatives there
directly by prime valuations.

\textbf{Solution 6.} Write \(\beta_j=p_j/q_j\) and \(\beta_n=p_n/q_n\)
in lowest terms, with positive denominators and \(p_n\ne0\). All
numerator and denominator absolute values are at most \(B\). The
unreduced fraction for \(-\beta_j/\beta_n\) has numerator \(-p_jq_n\)
and denominator \(q_jp_n\), each of absolute value at most \(B^2\).
Reduction cannot increase either, proving the height bound. Also
\(|\beta_n|=|p_n|/q_n\ge1/B\). Therefore \[
|\Lambda|=|\beta_n||\Lambda_*|
>B^{-1}(B^2)^{-C_*}=B^{-(1+2C_*)}.
\] At \(d=1\), Lemma 4.8 has \(r=4\), \(K_d=2\), and exponent
\(1+5C_*\). The sharper rational calculation avoids the deliberately
coarse general compositum and conjugate bounds. Both estimates are
uniform in the same variable \(B\).

\subsection{Prerequisites and
continuation}\label{prerequisites-and-continuation}

Theorem 4.1 is stated here and fully proved in
\href{../TR-BAKER-05.html}{\emph{Effective lower bounds II: proof of
Baker's theorem}}. Every arithmetic tool used to prepare that proof has
been proved above. The existing \emph{Number fields} lesson supplies the
elementary algebraic-integer criterion at exactly the required
generality. The later many-logarithm lesson owns the modern quantitative
refinements; the brief historical discussion above supplies neither a
substitute proof nor an additional prerequisite for the next lesson.

\subsection{References}\label{references}

\begin{itemize}
\tightlist
\item
  {[}Waldschmidt 2000{]} Michel Waldschmidt, \emph{Diophantine
  Approximation on Linear Algebraic Groups}, Springer, 2000.
  \href{https://webusers.imj-prg.fr/~michel.waldschmidt/articles/pdf/dalag.pdf}{Author's
  full text}. Effective homogeneous and nonhomogeneous estimates,
  integer-valued polynomial bases and Matveev's derivative-denominator
  method.
\item
  {[}Waldschmidt 2003{]} Michel Waldschmidt, ``Linear Independence
  Measures for Logarithms of Algebraic Numbers,'' in \emph{Diophantine
  Approximation}, Lecture Notes in Mathematics 1819, Springer, 2003,
  249--344.
  \href{https://webusers.imj-prg.fr/~michel.waldschmidt/articles/pdf/SpringerLN1819-2003.pdf}{Author's
  full text}. Baker's auxiliary-function method, normalized derivations
  and quantitative extrapolation.
\end{itemize}

\subsection{Editable source}\label{editable-source}

\href{TR-BAKER-04.md}{Markdown source}.

\end{document}
